% Solución explícita y clasificación del universo descendiente.

\section{Solución regular explícita de la continuación cosmológica descendiente}
\label{sec:ivv-descendant-explicit}

Considérese la ecuación efectiva homogénea de Einstein--Cartan para
\(n=3\),
\begin{equation}
 H^2=C\bigl(Aa^{-3}-Ba^{-6}\bigr),
 \qquad A,B,C>0.
 \label{eq:ivv-friedmann-descendant}
\end{equation}
Dados \(t_b\in\mathbb R\), defínase
\begin{equation}
 a(t)^3=\frac BA+\frac{9AC}{4}(t-t_b)^2.
 \label{eq:ivv-scale-factor-descendant}
\end{equation}

\begin{theorem}[Existencia no vacía y completitud de la continuación descendiente]
\label{thm:ivv-descendant-existence}
La función \eqref{eq:ivv-scale-factor-descendant} es una solución global de
\eqref{eq:ivv-friedmann-descendant}. Posee un único mínimo
\(a_b=(B/A)^{1/3}>0\), satisface
\[
 H(t_b)=0,
 \qquad
 \dot H(t_b)=\frac{3A^2C}{2B}>0,
\]
y define sobre \(\mathbb R\times\mathbb T^3\) una métrica FLRW lisa,
causalmente orientable, con invariantes de curvatura acotados y geodésicas
causales completas. Las regiones \(t<t_b\) y \(t>t_b\) realizan la rama
parental contractiva y la continuación descendiente expansiva, unidas por la
hipersuperficie regular \(t=t_b\).
\end{theorem}

\begin{proof}
Poniendo \(y=a^3\), se obtiene
\[
 H=\frac{\dot y}{3y}=\frac{3AC(t-t_b)}{2y},
 \qquad
 C(Aa^{-3}-Ba^{-6})
 =C\frac{Ay-B}{y^2}
 =\frac{9A^2C^2(t-t_b)^2}{4y^2}=H^2.
\]
La positividad de \(B/A\) excluye \(a=0\). Tanto \(H\) como \(\dot H\)
son acotados; los invariantes polinómicos de curvatura de la métrica FLRW
plana también lo son. El tiempo cósmico recorre \(\mathbb R\), y el parámetro
afín de las geodésicas nulas contiene \(\int a(t)\,dt\), que diverge en ambos
extremos porque \(a(t)\asymp |t|^{2/3}\). La longitud propia temporal diverge
asimismo. No existe, por tanto, borde causal finito ni singularidad en el
rebote.
\end{proof}

\begin{theorem}[Obstrucción del cobordismo ramificado clásico]
\label{thm:ivv-cobordism-obstruction}
En la clase de variedades lorentzianas lisas, no degeneradas, orientables en
el tiempo y globalmente hiperbólicas, una evolución entre superficies de
Cauchy compactas no puede producir una ramificación con cambio de topología.
Toda foliación de Cauchy es difeomorfa a \(\mathbb R\times\Sigma\); por ello,
las hojas inicial y final son difeomorfas. El cobordismo ramificado queda
obstruido dentro de esta clase.
\end{theorem}

Así quedan resueltos los \(2\) objetos geométricos: la continuación mínima
posee la solución completa anterior, y
la ramificación clásica con cambio topológico está excluida por el teorema de
obstrucción. Una geometría degenerada o una teoría cuántica de cambio
topológico constituye otro dominio y no modifica estos resultados.
